%!TEX root = parita-msc.tex

\begin{doublespace}
\begin{center}
\textbf{Abstract}
\end{center}

The \ac{www}, 
as defined by the \ac{w3c}, 
is ``the universe of network-accessible information, 
the embodiment of human knowledge''\footnote{\tt https://www.w3.org/WWW/}. 
With development over the years, 
the global web of linked documents has become the global web of linked data, 
which is popularly known as the Semantic Web.

The Semantic Web is now an established topic of research but there remain many unanswered questions. 
\ac{lod}, 
as the highest standard for linked data, 
is a worthy goal for researchers and data providers who seek to contribute to the development of the Semantic Web. 
Data that are linked and open facilitate discovery and reuse. 
Likewise, vocabularies used to describe the semantic relationships amongst data
that are linked and open facilitate discovery and reuse. 
Yet, it is not always a straightforward matter to take data that may be somehow accessible on a webpage written in \ac{html} and publish it as \ac{lod}.
Many complimentary technologies are used in handling \ac{lod} on the web, 
particularly the \ac{rdf}, 
the \ac{owl}, 
and the \ac{sparql}. 

Vandenbussche et al.~\cite{pierre2017vocabularies} in 2017 stated that 
\say{the \ac{lov} initiative gathers and makes visible indicators such as the interconnections between vocabularies and each vocabulary's version history, 
along with past and current editor (individual or organization)}. 
This thesis work is in contribution to the same initiative and to Berners-Lee's vision. 
As a part of this research, 
a model for survey ontology was developed in order to assist the study which is described in detail in Chapter~\ref{chap:proofOfConcept}.

\end{doublespace}